\documentclass[a4paper]{article}
\usepackage{tikz,cancel}
\usepackage[margin=1cm]{geometry}
\usepackage[spanish]{babel}
\title{"Impossible" question stumps students}
\author{Ferreyra Gustavo}
\begin{document}
\maketitle
\begin{center}
    \begin{tikzpicture}[scale=0.5]
\fill[cyan] (0,0) circle (4);
\draw[fill=white] (4,0) circle (4);
\draw[fill=white] (-4,0) circle (4);
\draw (0,0) circle (4);
\draw (0,0) circle (2pt);
\draw (-4,0) circle (2pt);
\draw (4,0) circle (2pt);
\draw (0,-5) node {r=4};
\end{tikzpicture}
\end{center}

\begin{center}
\begin{tikzpicture}[scale=0.5]
\draw[fill=white] (4,4) arc (90:270:4);
\draw[fill=white] (-4,-4) arc (-90:90:4);
\draw (0,0) circle (4);
\draw (0,0) circle (2pt);
\draw (-4,0) circle (2pt);
\draw (4,0) circle (2pt);
\draw[dashed] (-4,-4)rectangle(4,4);
\end{tikzpicture}
\end{center}
Se puede calcular calculando el área total y luego restando 2 circulos y el área de contra la recta. El área contra la recta es simetrica y se repite 8 veces.

$A=4r^2-\pi r^2 - 8A_R=$
\begin{center}
\begin{tikzpicture}
\draw[dashed,fill=cyan] (3.45,0) rectangle (4,2);
\draw[fill=white] (4,0) arc (0:30:4)-cycle;
\draw (0,0) -- node[above]{r}( 4,0);
\draw (0,0) --(30:4);
\end{tikzpicture}
\end{center}
Para calcular el área simétrica se suman el área del triángulo interno y el área del rectángulo, luego se le resta el área del sector.

$A_\Delta=\frac{r \cdot \sen(\pi/3) \times r\cdot \cos(\pi/3)}{2}=$

$A_{sector}=\pi/6 \cdot r^2$

$A_{rect}= r \cdot \cos(\pi/3)\times (r-r\cdot \sen(\pi/3)) $

como $\cos(\pi/3)=1/2$ y $\sen(\pi/3)=\sqrt{3}/2$

$A_\Delta=\frac{r \cdot 1/2 \times r\cdot \sqrt{3}/2}{2}=(r^2\sqrt{3}/8)$

$A_{sector}=\pi/12 \cdot r^2$

$A_{rect}=  r \cdot \cos(\pi/3)\times (r-r\cdot \sen(\pi/3))=$

$=r^2( \cos(\pi/3)\times (1-1\cdot \sen(\pi/3))=$

$=r^2( \frac{1}{2}\times (1-\sqrt{3}/2))=$

$=r^2(\frac{2-\sqrt{3}}{4}) $

$A_R=r^2\sqrt{3}/8+ r^2 (\frac{2-\sqrt{3}}{4})-\pi/12 \cdot r^2=r^2(\sqrt{3}/8+ (\frac{4-2\sqrt{3}}{8})-\pi/12)=$ 

$=r^2(\cancel{\sqrt{3}/8}+ (\frac{1}{2}\cancel{-2}\frac{\sqrt{3}}{8})-\pi/6)=r^2 (1/2-\sqrt{3}/8-\pi/12)$

$=r^2 (6/12-\sqrt{3}/8-\pi/12)=r^2 \frac{6-3/2\sqrt{3}-\pi}{12}$

Resultando entonces 

$A=r^2(4-\pi)-8 A_R=r^2(4-\pi)-8r^2 \frac{6-3/2\sqrt{3}-\pi}{12}=$

$r^2(4-\pi-8\frac{6-3/2\sqrt{3}-\pi}{12}))r^2(4-\pi-2\frac{6-3/2\sqrt{3}-\pi}{3}))$

$=r^2(\cancel{4}-\pi \cancel{-4}+\sqrt{3}+2\pi/3)=r^2 (\sqrt{3}-\pi/3)=16 (\sqrt{3}-\pi/3)\approx 10,96$

%\section{verification}

%$A<(4 \times 4=16)$

%$A<(pi/3 \approx 16.74)$


\end{document}